\documentclass[tikz,border=10pt]{standalone}
\usepackage{tikz}
\usepackage{amsmath}
\usepackage{amssymb}
\usepackage{ctex}
\usetikzlibrary{calc, positioning, arrows.meta, decorations.pathreplacing}

% A = [[3,1],[0,2]]
% 特征值 lambda1=3 (v1=(1,0)), lambda2=2 (v2=(1,-1))
% 普通向量: va=(1,2)->Ava=(5,4), vb=(-1,1)->Avb=(-2,2), vc=(0,1)->Avc=(1,2)

\begin{document}
\begin{tikzpicture}[>=Stealth, scale=1.15]

%% ====== 左子图：普通向量（方向改变）======
\begin{scope}[xshift=0cm]

  % 网格与坐标轴
  \draw[gray!20, thin] (-2.3,-2.3) grid (5.3,5.3);
  \draw[->,gray!50, thin] (-2.4,0) -- (5.5,0) node[right,font=\scriptsize] {$x$};
  \draw[->,gray!50, thin] (0,-2.4) -- (0,5.5) node[above,font=\scriptsize] {$y$};
  \fill (0,0) circle (1.5pt) node[below left,font=\scriptsize] {$O$};

  % --- 普通向量 va=(1,2) 蓝色实线 ---
  \draw[->,thick,blue!70!black,>=Stealth]
    (0,0) -- (1,2)
    node[left,font=\small,blue!70!black] {$\mathbf{v}_a$};
  % Ava = (3*1+1*2, 0*1+2*2) = (5,4) — 红虚线
  \draw[->,thick,red!70!black,dashed,>=Stealth]
    (0,0) -- (5,4)
    node[right,font=\small,red!70!black] {$A\mathbf{v}_a$};
  % 方向改变提示弧
  \draw[->,gray!50,thin] (0.6,1.2) to[bend left=18] (2.5,2.0);

  % --- 普通向量 vb=(-1,1) 蓝色实线 ---
  \draw[->,thick,blue!50!black,>=Stealth]
    (0,0) -- (-1,1)
    node[above left,font=\small,blue!50!black] {$\mathbf{v}_b$};
  % Avb = (3*(-1)+1*1, 0*(-1)+2*1) = (-2,2) — 红虚线
  \draw[->,thick,red!50!black,dashed,>=Stealth]
    (0,0) -- (-2,2)
    node[above left,font=\small,red!50!black] {$A\mathbf{v}_b$};

  % --- 普通向量 vc=(2,1) 蓝色实线 ---
  \draw[->,thick,blue!90!teal,>=Stealth]
    (0,0) -- (2,1)
    node[below right,font=\small,blue!90!teal] {$\mathbf{v}_c$};
  % Avc = (3*2+1*1, 0*2+2*1) = (7,2) — clip到 (5,2)缩放显示: 使用(5,1.43)
  % 实际: (7,2), 太大，缩短为0.7倍=(4.9,1.4)
  \draw[->,thick,red!90!orange,dashed,>=Stealth]
    (0,0) -- (4.9,1.4)
    node[right,font=\small,red!90!orange] {$A\mathbf{v}_c$};

  % 子标题
  \node[font=\small\bfseries,align=center] at (1.5,-2.0)
    {普通向量：方向改变};

\end{scope}

%% ====== 右子图：特征向量（只拉伸，方向不变）======
\begin{scope}[xshift=9cm]

  % 网格与坐标轴
  \draw[gray!20, thin] (-2.3,-3.3) grid (5.3,3.3);
  \draw[->,gray!50, thin] (-2.4,0) -- (5.5,0) node[right,font=\scriptsize] {$x$};
  \draw[->,gray!50, thin] (0,-3.4) -- (0,3.5) node[above,font=\scriptsize] {$y$};
  \fill (0,0) circle (1.5pt) node[below left,font=\scriptsize] {$O$};

  % --- 特征向量 v1=(1,0), lambda1=3 ---
  \draw[->,very thick,green!60!black,>=Stealth]
    (0,0) -- (1,0)
    node[below,font=\small,green!60!black] {$\mathbf{v}_1$};
  % A*v1 = (3,0) = 3*(1,0) 方向相同，长度3倍
  \draw[->,very thick,green!40!black,dashed,>=Stealth]
    (0,0) -- (3,0)
    node[below right,font=\small,green!40!black] {$A\mathbf{v}_1 = 3\mathbf{v}_1$};
  % 拉伸标注
  \draw[<->,orange,thick] (0,0.35) -- (3,0.35)
    node[midway,above,font=\scriptsize,orange] {$\lambda_1=3$};

  % --- 特征向量 v2=(1,-1), lambda2=2 ---
  % A*(1,-1) = (3+(-1), 0+(-2)) = (2,-2) = 2*(1,-1) ✓
  \draw[->,very thick,green!70!teal,>=Stealth]
    (0,0) -- (1,-1)
    node[right,font=\small,green!70!teal] {$\mathbf{v}_2$};
  \draw[->,very thick,green!50!teal,dashed,>=Stealth]
    (0,0) -- (2,-2)
    node[right,font=\small,green!50!teal] {$A\mathbf{v}_2 = 2\mathbf{v}_2$};
  % 拉伸标注
  \node[font=\scriptsize,orange,fill=white,inner sep=2pt] at (2.2,-0.5)
    {$\lambda_2=2$};

  % 说明文字
  \node[font=\small,fill=green!8,draw=green!50,rounded corners=3pt,
        inner sep=6pt,align=center]
    at (1.5, 2.8) {
      特征向量：方向不变\\
      $A\mathbf{v}=\lambda\mathbf{v}$，只伸缩
    };

  % 子标题
  \node[font=\small\bfseries,align=center] at (1.5,-3.0)
    {特征向量：方向不变，只缩放};

\end{scope}

%% ====== 矩阵标注框（两图之间）======
\node[font=\small, fill=white, draw=gray!50, thin, rounded corners=4pt,
      inner sep=8pt, align=center]
  at (4.5, -3.2) {
    $A=\begin{pmatrix}3&1\\0&2\end{pmatrix}$，
    特征值 $\lambda_1=3$，$\lambda_2=2$
  };

%% 整体标题
\node[font=\large\bfseries] at (4.5, 6.2)
  {特征向量的几何意义};

\end{tikzpicture}
\end{document}
